% subsection: 等差型 \, $a_{n+1}=a_n+d$

  \subsection{\texorpdfstring{等差型 \, $a_{n+1}=a_n+d$}{等差型}}\label{節：等差型}
  \begin{bookpoint}
    \textbf{差を追う}\par
    一回ごとに足す量が一定なら,初項に増分を足し戻す.項の番号と更新の回数の違いを確かめよう.
  \end{bookpoint}
    一歩ごとに$d$を足す.$n\ge3$で直前の項を順に代入すると,
    \begin{align*}
      a_n&=a_{n-1}+d=a_{n-2}+2d=\cdots\\
         &=a_1+(n-1)d
    \end{align*}
    となり\bookterm{general}{一般項}が示される.$n=1,2$では直接確かめられる.
    例えば$a_4=a_3+d=a_2+2d=a_1+3d$である.項は$n$個,更新は$n-1$回なのである.

    \noindent \bookblocklink{example-arithmetic-1}{problem}{answer}{【例題】}\par\nobreak
    \begin{enumerate}[label=(\arabic*),parsep=3pt,format=\bookitemlink{example-arithmetic-1}{problem}{answer}]
      \item $\displaystyle a_1=1,\, a_{n+1}=a_n+1$
      \item $\displaystyle a_1=\frac{1}{2},\, a_{n+1}=a_n+\frac{1}{4}$
      \item $\displaystyle a_1=-\sqrt{3},\, a_{n+1}=a_n-\sqrt{2}$
    \end{enumerate}

    \noindent\bookblocklink{example-arithmetic-1}{hint}{problem}{【方針】}\par\nobreak
    \begin{enumerate}[label=(\arabic*),parsep=3pt,format=\bookitemlink{example-arithmetic-1}{hint}{problem}]
      \item 等差型の\bookterm{general}{一般項}
            $\displaystyle a_n=a_1+(n-1)d$
            を用います.
      \item 公差は$\displaystyle d=\frac14$です.
            等差型の\bookterm{general}{一般項}に代入します.
      \item 公差は$\displaystyle d=-\sqrt2$です.
            負の公差の場合も同じ公式を使います.
    \end{enumerate}

    \noindent\bookblocklink{example-arithmetic-1}{answer}{problem}{【解答】}\par\nobreak
    \begin{enumerate}[label=(\arabic*),parsep=3pt,format=\bookitemlink{example-arithmetic-1}{answer}{problem}]
      \item $\displaystyle a_n=1+(n-1)=n$
      \item $\displaystyle a_n=\frac12+\frac{n-1}{4}
            =\frac{n+1}{4}$
      \item $\displaystyle a_n=-\sqrt3-(n-1)\sqrt2$
    \end{enumerate}

    \noindent\bookblocklink{example-arithmetic-1}{solution}{problem}{【解法】}\par\nobreak
    (1)は$a_1=d=1$なので$a_n=1+(n-1)=n$となる.初項は$1$,隣り合う項の差は$(n+1)-n=1$で,元の条件に合う.
    (2),(3)も初項と公差を代入する.
